\documentclass[12pt]{article}
\usepackage[ left=0.5in,right=0.5in ]{geometry}
\usepackage{../styles/ntalgo}
\usepackage[backend=biber,style=alphabetic]{biblatex}
\addbibresource{refs.bib}

\begin{document}

\title{Karger Stuff}
\author{Navid Tajkhorshid}

\maketitle

Notes based on various works of Karger and friends (mainly~\cite{BK02}) on
sparse graphs which preserve cuts. I don't really care about runtime here.

\section{Definitions}

We're working with an undirected, possibly weighted graph $G = (V, E)$. Let
$c : E \to \R_{\geq 0}$ denote a capacity function on the edges; if $G$ is
unweighted, $c \equiv 1$. A \emph{cut} is a set of vertices
$\emptyset \neq U \subsetneq V$. The capacity of this cut is
$c(\cut{U}) \defeq \sum_{e \in \cut{U}}c(e)$. I'll also use cut to refer to the
set of edges leaving $U$, when the context is clear. Generally, we're interested
in the values of various cuts. We start by looking at the minimum cut.

\section{Karger's Algorithm}

This is a very slick algorithm to compute the minimum cut in $G$. While $G$ has
more than 2 vertices, we pick a random edge $e = (u, v)$ proportional to its
capacity (that is, the chance of picking $e$ is $\frac{c( e )}{C(E)}$). We then
contract this edge: merge the vertices $u$ and $v$ into a new vertex $w$, with
edges $(w, x)$ for all original edges $(u, x)$ and $(v, x)$ (note that the graph
may become a multigraph). When we get down to two vertices $u$ and $v$, output
$E(\Set{u}, \Set{v})$.

To argue that the outputted $S$ is a good cut, we'll fix some minimum cut $S^*$
and analyze its chance of survival throughout the algorithm (that is, the chance
that no edge in $S^*$ is contracted). Let $\lambda^* = c(S^*)$. The chance we
contract an edge in $S^*$ in a given step is the $c(S^*)/c(E)$ (note that $E$
represents the remaining edges after some steps of contraction). The total
weight of all edges is $c(E)$, so the average weight of the singleton cuts
$\Set{v}$ for each $v \in V$ is $2c(E)/n$ (each edge is counted twice). This
means some singleton has weight at most $2c(E)/n$, which cannot be smaller than
$\lambda^*$. The probability of failure is then at most
\[
  \frac{2c(E)/n}{c(E)} = \frac{2}{n}.
\]
The chance of survival is $1 - \frac{2}{n} = \frac{n - 2}{n}$, so over all steps
we have
\begin{align*}
  \Pr{S^* \text{ survives}} \geq \prac{n - 2}{n}\prac{n - 3}{n - 1}\prac{n - 4}{n - 2}\cdots\prac{2}{4}\prac{1}{3}
                            =    \frac{2}{n(n - 1)}
                            =    1\bigg/\binom{n}{2}.
\end{align*}

Naively, one can repeat the algorithm $\bigO{n^2\log n}$ times and return the
best cut with a high probability guarantee. There is also more one can do to
achieve better running time, but we're not interested in that right now.

\subsection{Bounding Cuts}

An interesting feature of the argument above is that it works for any specified
min cut $S^*$. Since the algorithm returns any min cut with probability at least
$\binom{n}{2}^{-1}$, there can be at most $\binom{n}{2}$ distinct min cuts in
$G$. Note that this is achieved by the graph $C_n$--any pair of edges
disconnects the graph.

It will be useful to generalize this bound to the number of
\emph{$\alpha$-approximate} min cut (a min cut of value at most
$\alpha\lambda^*$). By rounding $\alpha$ up, assume it's an integer. If we try
to step through the same analysis--this time with a failure probability per step
bounded by $\frac{2\alpha}{n}$--we run into an issue when $n$ gets too small.

We can fix this via early termination: When $n \leq 2\alpha$ stop the algorithm.
Note for $\alpha = 1$ this is the same as $n \leq 2$, i.e.\ it's identical to
the standard algorithm. The probability that an $\alpha$-approximate min cut
survived is at least
\begin{align*}
  \prac{n - 2\alpha}{n}\prac{n - 2\alpha - 1}{n - 1}\prac{n - 2\alpha - 2}{n - 2}\cdots\prac{2}{2\alpha + 2}\prac{1}{2\alpha + 1}
  = \frac{(2\alpha)(2\alpha - 1)\cdots(2)(1)}{n(n - 1)\cdots(n - 2\alpha + 1)}
  = 1\bigg/\binom{n}{2\alpha}.
\end{align*}
Note however, that since we terminated when $n = 2\alpha$, there can be many
cuts that survived. In particular, there are $2^{2\alpha-1} - 1$ such cuts (for
each cut, count the smaller, nonempty side). There are thus at most
$\binom{n}{2\alpha}2^{2\alpha-1}$ approximate minimum cuts. Since
$(2\alpha)! \leq 2^{2\alpha-1}$, we have
\begin{align*}
  \binom{n}{2\alpha}2^{2\alpha-1} \leq n^{2\alpha}
\end{align*}
as an upper bound.

\section{Preserving Cuts}

\def\U{\mathcal{U}}

Suppose we want to approximate find the min-cut in a graph, and say $G$ is
dense. It would be nice to operate on a sparser graph. In particular, say we
could compute a graph $G' = (V, E, c')$ such that every cut $U \subseteq V$ in
$G'$ obeyed
\[
  (1 - \eps)c(\cut{U}_G) \leq c'(\cut{U}_{G'}) \leq (1 + \eps)c(\cut{U}_G),
\]
that is, the cuts in $G'$ approximate those in $G$. Then computing an
approximate min cut in $G'$ would give an approximate min cut in $G$. A simple
idea is to sample down the number of edges and scale capacities up: pick some
probability $p$, include $e \in E$ into $G'$ with probability $p$ and set
$c'(e) \defeq c(e)/p$. It's not hard to see that
$\E{c'(\cut{U}_{G'})} = c(\cut{U}_G)$ for every $U \subseteq V$.

Expectation isn't good enough, of course. Let's fix a cut $U \subseteq V$ and
see how concentrated $c(\cut{U}_{G'})$ is about its expectation. We want a
multiplicative bound, so we don't have to worry about the scaling. Also scale
capacities so that $c(e) \leq 1$ for each edge. Then,
$c(\cut{U}_{G'}) = \sum_{e \in \cut{U}_G}c(e)X_e$, where $X_e$ is the indicator
for $e$ being sampled. Let $X$ denote this random sum--$\E{X} = pc(\cut{U}_G)$.
By a standard Chernoff inequality, we have
\[
  \Pr{X \notin (1 \pm \eps)\E{X}} \leq 2e^{-\E{X}\eps^2/3}
                                  =    2e^{-pc(\cut{U}_G)\eps^2/3}.
\]
To get a high probability guarantee, we'd want $p \geq \bigOmega{\frac{\log
    n}{\eps^2 c(\cut{U}_G)}}$ for every cut $U$. So let's set $p = \bigTheta{\frac{\log n}{\eps^2 \lambda}}$, where $\lambda$ is some estimate of
the minimum cut value which we assume is accurate up to constant factors. Then
we get that a given cut is preserved with high probability. Unfortunately, there
are $\bigTheta{2^n}$ cuts, so we can't do a naive union bound. There is hope
though: from the discussion above, the number of cuts with small values is
polynomial. The larger the value of a cut is, the more aggressive a bound we get
from the analysis. To formalize this, let $\U_\alpha$ denote the set of cuts
whose value is between $\alpha\lambda^*$ and $2\alpha\lambda^*$. By the bound
above, $\sizeof{\U_\alpha} \leq \bigO{n^{4\alpha}}$, since each cut is
$2\alpha$-approximate. Now consider a cut $U \in \U_\alpha$. Running through the
probability bounds again, we get
\[
  \Pr{\text{$U$ isn't adequately preserved}}
  \leq 2e^{-\bigOmega{pc(\cut{U}_G)\eps^2}}
  =    2e^{-\bigOmega{\prac{\log n}{\eps^2\lambda^*}\alpha\lambda^*\eps^2}}
  =    2n^{-\bigOmega{\alpha}}.
\]
Then, by a union bound over all $\bigO{n^{4\alpha}}$ elements, we have a high
probability guarantee for $\U_\alpha$ being preserved. We can do this for all
$\alpha$ up through $n$. For the remaining cuts of size larger than
$n\lambda^*$, each is preserved with probability $1 - n^{-\bigOmega{n}}$, so a
bound over all $2^n$ cuts suffices.

\subsection{Better Bounds?}

In the remainder of this section, assume $G$ is unweighted. The obvious
weaknesses of the above approach are two-fold. First, it requires an estimate on
the value of the minimum cut in $G$. Moreover, if that value $\lambda^*$ is
small, we barely get any sparsification: if $\lambda^*\eps^2 = O(1)$, then a
constant factor of the edges are preserved in expectation.

The following example illustrates the algorithm's inefficiency. Suppose $G$ is
the unweighted union of two copies of $K_n$, joined by a single edge (a sort of
``dumbbell graph''). The min cut value in $G$ is $1$, so the sparsification
algorithm would keep every edge ($\frac{\log n}{\eps^2} \geq 1$). This is quite
wasteful: a less oblivious algorithm could separately sparsify the two cliques,
which have minimum cut value $n - 1 \approx n$, and keep the edge between them
with probability 1. This leads to
$\bigO{n^2\log n/(n\eps^2)} = \bigO{n\log n /\eps^2}$ remaining edges! Note that
we have to worry about scaling now that sampling probabilities are
nonuniform--sampled edges within the cliques get weight $n\eps^2/\log n$ while
the bridge gets weight $1$.

\begin{center}
  [Insert figure of dumbell graph turning into sparsified one]
\end{center}
\begin{xydiagram}
  K_n \ar@{-}[rr]^{\text{1}} &  & K_n
\end{xydiagram}

Every cut in $G$ decomposes into cuts in the cliques (with are preserved whp),
and perhaps the single nonempty cut along the ``bridge'' edge. The bridge edge
is preserved, and cuts within the cliques are preserved to $1 \pm \eps$, so all
cuts are adequately preserved.

Let's try a harder case: Suppose now $G$ consists of two cliques joined by
$\bigTheta{\sqrt{n}}$ bridge edges. Of course, we can sample down the cliques
like before and keep every bridge, but in some soft sense this is wasteful, and
it's not hard to come up with examples where this doesn't work\footnote{Indeed,
we can include a dense graph with connectivity $\sqrt{n}$ in the middle of the
bridges; including all its edges would kill any asymptotic sparsification.}.
Really, we'd like to sample down the bridge edges as well: if a sampling
probability of $\frac{\log n}{n\eps^2}$ works for the cliques of connectivity
$\bigTheta{n}$, maybe we can be brave and sample down the bridge edges with
connectivity $\bigTheta{\sqrt{n}}$ with probability
$\frac{\log n}{\sqrt{n}\eps^2}$. Again, each edge, if sampled, is included in
$G'$ with capacity reciprocal to its sampling probability.

Arguing that cuts are preserved becomes much trickier. The main annoyance is
that cuts that involve bridge edges can be polluted by clique edges.
Concentration bounds can't work since sampled clique edges have large weight.
The trick is the following. We'll write $G$ as the weighted sum of three
capacitated graphs $G_1$, $G_2$, and $G_3$. In $G_1$, we'll include all edges,
but set the weight of clique edges to $1/\sqrt{n}$. $G_2$ and $G_3$ are copies
of the cliques. We sample an edge with probability $p_e$ still; the downsampled
$G_i$s ($G_i'$) contain $e$ with $1/p_e$ times its weight in each. We want to
show that cuts are preserved within a $(1 \pm \eps)$ factor in each of these
graphs; if we have this, then cuts are also preserved in
\[
  G = G_1 + \paren{1 - 1/\sqrt{n}}(G_2 + G_3)
\]
(exercise). This is clear for $G_2$ and $G_3$, since their min cut is
$\bigTheta{n}$.

The minimum cut in $G_1$ (like in $G$) has value $\bigTheta{\sqrt{n}}$: if any
clique is cut we it automatically captures at least $(n - 1)/\sqrt{n}$ capacity;
otherwise it cuts all $\sqrt{n}$ unit-weight bridge edges. A sampled edge
contributes
\[
  \bigO{\frac{\sqrt{n}\eps^2}{\log n}} = \bigO{\frac{\lambda_1^*\eps^2}{\log n} }
\]
capacity (clique edges were scaled down for this reason), that is, each
contributes a $\frac{\eps^2}{\log n}$ fraction of the min cut. From the previous
section, we know that all cuts in $G_1'$ are within a $(1 \pm \eps)$ factor of
those in $G_1$, and we're done.

\begin{center}
  [Figure: $G$ as a conic combination of $G_1$, $G_2$, and $G_3$,
    where each preserve cuts.]
\end{center}
This suggests a more general algorithm that looks for subgraphs with high
connectivity and samples down their edges much more aggressively\footnote{In
fact, this is implicitly done in the special case of connected components--it's
easy to see why we can assume WLOG that $G$ is connected.}. We next build out
some definitions that let us formalize this.

\subsection{Strong $\kappa$-Connectivity}

Recall some definitions: a graph $G$ is $\kappa$-connected (or has connectivity
$\kappa$) the min cut is at least $\kappa$. Two vertices $u$ and $v$ are
$\kappa$ connected is the min $(u, v)$-cut has value at least $\kappa$. A
stronger notion is the following.
\begin{definition}
  A pair of vertices $u$ and $v$ (or an edge $(u, v)$) is
  \emph{$\kappa$-strongly connected} if there is some $U \subseteq V$ containing
  $u$ and $v$ such that $G[U]$ has connectivity at least $\kappa$.
\end{definition}

Observe that $\kappa$-strong connectivity is an equivalence relation. The only
interesting part is transitivity, so consider $u, v, w$ where $u, v$ and $v, w$
are $\kappa$-strongly connected. We then have sets $U_1 \ni u, v$ and
$U_2 \ni v, w$ such that the corresponding subgraphs are $\kappa$-connected. We
claim that $G[ U_1 \cup U_2 ]$ is also $\kappa$-connected. The reason is that
$U_1$ and $U_2$ have nonempty intersection, so any cut of their union must cut
one of the subsets.

There is another more algorithmic way to see the equivalence relation, where we
build the partition into equivalence classes:
% badness-format skip
\begin{centeralgo}
  \textsc{Split}($G, \kappa$):\+
  \\ If $G$ has connectivity $\kappa$, return $\Set{V}$
  \\ Otherwise, let $U$ be a cut of value less than $\kappa$
  \\ return $\textsc{Split}(G[U], \kappa) \cup \textsc{Split}(G[V \setminus U], \kappa)$
\end{centeralgo}

Correctness comes down to two observations. First, it is clear that no pair of
vertices separated by $U$ can be $\kappa$-strongly connected. Second, if, say
$u, v \in U$ are $\kappa$-strongly connected, then the any witness set $U'$
cannot cross $U$: if it did, then then $U \cap U'$ would be a cut with value
less than $\kappa$ in $G[U']$. We'll call the subgraphs induced by these
components \emph{$\kappa$-strongly connected components}.

The second observation also implies that the partition gets finer as $\kappa$
increases. For $\kappa_1 < \kappa_2$, any equivalence class under $\kappa_2$
must be contained in an equivalence class under $\kappa_1$.

A useful way to think about these is as a sequence of finer and finer
partitions.
\def\P{\mathcal{P}}
\begin{exercise}
  % TODO: This is technically ill-formed
  Modify the algorithm above to output a sequence of partitions
  $\P_0, \P_1, \dots, \P_{n-1}$, where $\P_i$ induces the set of $i$-strongly
  connected components in $G$,
\end{exercise}

\begin{center}
  Figure: A graph getting decomposed into stronger and stronger components. At
  least 3 levels of decomposition.
\end{center}

\subsection{Ideal Sparsification}

Suppose we are given the edge strengths $\kappa_e$ for each $e \in E$. Let's try
to extend the decomposition technique from the dumbbell graphs.

For each $e \in E$, we include it in $E'$ with probability
$p_e = \bigTheta{\dfrac{\log n}{\kappa_e\eps^2}}$; its weight if sampled becomes
$1/p_e$.\footnote{We technically cap $p_e$ at $1$, but we ignore this to clean
things up.}

We consider the following collection of graphs: for $\kappa < n$, let $G_\kappa$
denote $G$ restricted to edges of strength at least $\kappa$. If an edge is
included in $G_\kappa$, it is given capacity $1/\kappa_e$. Let $G_\kappa'$
denote $G_\kappa$ restricted to sampled edges; each has capacity
$c_\kappa(e)/p_e$. Now fix $\kappa$. We want to argue that $G_\kappa'$ preserves
cuts. It sufficed to argue this for a connected component $C$ of $G_\kappa$.

We claim that $C$ has minimum cut at least $1$. For any cut $U \subseteq V$, we
have
\begin{align*}
  c_{G_\kappa}(\cut{U}) =    \sum_{e \in \cut{U}}\frac{c(e)}{\kappa_e}
                        \geq \sum_{e \in \cut{U}} \frac{c(e)}{c(\cut{U})}
                        =    1,
\end{align*}
where the inequality is because the cut $U$ certifies that
$\kappa_e \leq c(\cut{U})$ for each $e \in \cut{U}$. (Why? This is a little
subtle).

If an edge is in $G_\kappa'$, it has weight
\[
  c_\kappa(e)/p_e = \bigO{\eps^2/\log n}.
\]
That is, the edge contributes $\eps^2/\log n$ fraction of the min cut. This
immediately implies that $G_\kappa'$ preserves cuts in $G_\kappa$ to a
$(1 \pm \eps)$ factor.

We can write $G$ as
\[
  G = \sum_{\kappa=1}^{n-1} G_\kappa,
\]
and by a union bound over all $n$ subgraphs, we preserve all cuts with high
probability.

\subsection{Number of Edges}

Of course, we still need to show that $G'$ is a small number of edges. Since an
edge $e$ is sampled with probability $\log n / (\kappa_e \eps^2)$, we have
\[
  \E{\sizeof{E'}} = \sum_{e \in E} \frac{\log n}{\kappa_e \eps^2}
                  = \prac{\log
    n}{\eps^2}\sum_{e \in E} 1/\kappa_e.
\]

We can bound the last term with the following induction. We claim that
$\sum_{e \in E} 1/\kappa_e$ is bounded by $n - 1$. In the base case with a
single vertex, this is true ($\kappa_e$ is technically infinite). Otherwise
consider a min-cut $U$ of $G$. The edges in $\cut{U}$ all have strength at least
$\sizeof{\cut{U}}$ (since $\sizeof{\cut{U}}$ is the connectivity of $G$). Then,
\begin{align*}
  \sum_{e \in E}1/\kappa_e & = \sum_{e \in \cut{U}}1/\kappa_e + \sum_{e \in G[U]} 1/\kappa_e + \sum_{e \in G[V \setminus U]} 1/\kappa_e                      \\
                           & \leq 1 + \sizeof{U} - 1 + \sizeof{V \setminus U} - 1                                                       & \text{(induction)} \\
                           & = n - 1.
\end{align*}
Note that the induction works because, while edge strengths may decrease in
subgraphs, the bounds on edge strengths is based on the connectivity of the
subgraph. (FIX)

We are left with a sparsification algorithm that gives $\bigO{n \log n /
  \eps^2}$ edges in expectation (and with high probability, why?).

\section{Generalizing to Weighted Graphs}

Generalizing is mainly busywork\dots Essentially replace $1$ with $c(e)$ where
appropriate.
\printbibliography

\end{document}
